\documentclass{beamer}
\usepackage[utf8]{inputenc}

\usetheme{Berkeley}
\usecolortheme{default}

\setbeamertemplate{footline}[frame number]
\beamertemplatenavigationsymbolsempty


\title{Meta-Structure}
\author[]{IB 514: Scientific Writing}
\date[]{}

\AtBeginSection[]
{
  \begin{frame}
    \frametitle{Overview}
    \tableofcontents[currentsection]
  \end{frame}
}

\begin{document}

% -------------------------------------------------
\begin{frame}
  \titlepage
\end{frame}

\section{Audience and Journal Selection}

\begin{frame}
  \frametitle{Audience Identification}
  \begin{itemize}
    \item Know your target audience:
    \begin{itemize}
      \item Specialists in your field
      \item Interdisciplinary researchers
      \item Practitioners or policymakers
      \item General scientific audience
    \end{itemize}
    \item Audience shapes decisions about:
    \begin{itemize}
      \item Terminology and background explanation
      \item Level of technical detail
      \item Context and significance framing
      \item Visual presentation choices
    \end{itemize}
    \item Tailor writing to match audience expertise and interests
    \item Consider what your audience needs to know versus already knows
  \end{itemize}
\end{frame}

\begin{frame}
  \frametitle{Journal Selection}
  \begin{itemize}
    \item Identify appropriate target journals early
    \item Key considerations:
    \begin{itemize}
      \item Scope and readership of the journal
      \item Impact factor and citation metrics
      \item Publication timeline and backlog
      \item Author guidelines and manuscript format
      \item Open access policies and costs
    \end{itemize}
    \item Align your manuscript structure with journal requirements
    \item Review recent articles in target journal to understand expectations
    \item Consider journal fit before and during writing
  \end{itemize}
\end{frame}

\section{Methods and Materials}

\begin{frame}
  \frametitle{Methods Section Detail}
  \begin{itemize}
    \item Provide sufficient detail for reproducibility
    \item Include:
    \begin{itemize}
      \item Study design and justification
      \item Sampling methods and sample sizes
      \item Data collection protocols
      \item Equipment and software specifications (versions, parameters)
      \item Statistical methods and analysis procedures
      \item Quality control and validation measures
    \end{itemize}
    \item Balance completeness with readability
    \item Use subheadings to organize complex sections
    \item Reference protocols: cite published methods or provide in supplementary materials
    \item Specify materials: chemicals, organisms, equipment with model numbers
  \end{itemize}
\end{frame}

\begin{frame}
  \frametitle{Supplementary Materials}
  \begin{itemize}
    \item Use supplementary materials strategically
    \item Include in supplements:
    \begin{itemize}
      \item Detailed statistical analyses and raw data
      \item Extended methods or alternative approaches
      \item Additional figures and tables
      \item Raw data files (when permitted)
      \item Supplementary text and background information
    \end{itemize}
    \item Keep main text focused on key findings
    \item Reference supplementary materials clearly in main text
    \item Ensure main paper is complete and understandable independently
    \item Check journal policies for supplementary material format and limits
  \end{itemize}
\end{frame}

\section{Technical Elements}

\begin{frame}
  \frametitle{Integration of Equations}
  \begin{itemize}
    \item Introduce equations before presenting them
    \item Use equations to express complex relationships clearly
    \item Number important equations and reference them
    \item Guidelines:
    \begin{itemize}
      \item Define all variables and parameters
      \item Explain the logic and significance of equations
      \item Keep equations in the main text only if essential
      \item Consider moving complex equations to supplementary materials
    \end{itemize}
    \item Follow mathematical notation conventions
    \item Use consistent formatting and typography
    \item Explain derivations or cite sources for unfamiliar equations
  \end{itemize}
\end{frame}

\section{Title and Abstract}

\begin{frame}
  \frametitle{Title and Abstract Design}
  \begin{itemize}
    \item Title should be:
    \begin{itemize}
      \item Specific and descriptive (10-12 words optimal)
      \item Contain key terms for searchability
      \item Free of jargon when possible
      \item No abbreviations or citations
    \end{itemize}
    \item Abstract should:
    \begin{itemize}
      \item Summarize background, methods, results, and conclusions
      \item Be self-contained (readable without viewing paper)
      \item Meet journal length requirements
      \item Include primary findings and key quantitative results
      \item State implications and conclusions clearly
    \end{itemize}
    \item Both are the most-read parts of your paper
    \item Optimize for searchability and impact
  \end{itemize}
\end{frame}

\section{Organization and Flow}

\begin{frame}
  \frametitle{Signposting}
  \begin{itemize}
    \item Guide readers through your argument with effective signposting
    \item Use transitions and connectives:
    \begin{itemize}
      \item Topic sentences that signal direction
      \item Transitional phrases linking paragraphs and sections
      \item Explicit statements of relationships between ideas
    \end{itemize}
    \item Employ structural markers:
    \begin{itemize}
      \item Clear section headings and subheadings
      \item Numbered or bulleted lists for complex information
      \item Summary sentences to recap key points
    \end{itemize}
    \item Help readers understand:
    \begin{itemize}
      \item What comes next
      \item How ideas connect
      \item Why information matters
    \end{itemize}
  \end{itemize}
\end{frame}

\begin{frame}
  \frametitle{Structural Coherence}
  \begin{itemize}
    \item Maintain logical flow throughout the paper
    \item Structure provides the ``roadmap'' for readers
    \item Key aspects:
    \begin{itemize}
      \item Logical progression of ideas within sections
      \item Clear connections between sections
      \item Motivation established before presenting methods
      \item Results presented before interpretation
      \item Conclusions grounded in findings
    \end{itemize}
    \item Typical structure (IMRAD):
    \begin{itemize}
      \item Introduction: context and questions
      \item Methods: how you answered questions
      \item Results: what you found
      \item Discussion: what it means
    \end{itemize}
    \item Ensure each section builds on previous ones
  \end{itemize}
\end{frame}

\section{Learning from Examples}

\begin{frame}
  \frametitle{Analysis of Effective Papers}
  \begin{itemize}
    \item Learn by studying exemplary papers in your field
    \item Analyze structure and organization:
    \begin{itemize}
      \item How do authors present their problem?
      \item What background information do they provide?
      \item How do they organize results?
      \item How do they connect findings to broader context?
    \end{itemize}
    \item Examine writing strategies:
    \begin{itemize}
      \item Sentence construction and clarity
      \item Use of active vs.\ passive voice
      \item Paragraph organization and transitions
      \item Balance of technical detail and accessibility
    \end{itemize}
    \item Identify patterns in high-impact papers
    \item Model your own writing on effective examples
    \item Note what works and why
  \end{itemize}
\end{frame}

\end{document}
